\subsection{NN-G028: fixed bottom coefficients force an exact 3352-pair optimum}
\paragraph{Statement.}
Among all effective associative closed rank-three tables with
\[
 W=3D,\qquad U_1=3D,\qquad U_2=D,                    \tag{G028.1}
\]
the exact cancelling-pair minimum is $3352$.
Cell polynomials, contents and degrees are arbitrary;
$G_i/(W\delta_i)$ and $\Lambda_i$ may be arbitrary
rational functions of $D$. Only the actual table
coefficients are required to be integral, effective,
dominant and parity-correct. There are exactly two
minimizing parameter sets within this closed presentation.
One is (G027.4); the other replaces its $V_2$ by
$2D^2-D+4$ and adjusts $A_2$ accordingly.

The fixed coefficients (G028.1) are a substantive
restriction. This theorem closes a rational-function
optimization problem inside that restriction, and
makes no lower-bound assertion when $W$ or the two
$U_i$ vary.

\paragraph{A forced two-polynomial parameterization.}
The scalar identities used in NN-G011 are necessary
for associativity in this closed table: associativity
with $\Sigma$ gives
\[
 WS=G_1G_2,\qquad U_iS=G_i(G_i-D\delta_i),\qquad
 V_iS=G_i(DG_i-A_i),\quad A_i=\delta_i\Lambda_i.
\]
For example, $\Sigma(\gamma_1^1\gamma_1^2)$ gives
the first identity; multiplication by $\Sigma$ of
a same-cell two-letter product gives the second,
and of a three-letter reduction gives the third.
Thus no generic-embedding assumption is added in
the lower-bound argument.

Substituting (G028.1) yields
\[
 G_1-D\delta_1=G_2,\qquad
 3(G_2-D\delta_2)=G_1.
\]
Consequently
\[
 \begin{gathered}
 H=(\delta_1+\delta_2)/4,\qquad
 J=(\delta_1+3\delta_2)/2,\\
 G_1=6DH,\quad G_2=DJ,\quad S=2DHJ,\qquad
 \delta_1=6H-J,\quad\delta_2=J-2H.                 \end{gathered}\tag{G028.2}
\]
Both $H,J$ are effective integral symmetric Laurent
polynomials, hence belong to $\mathbb Z[D]$.
For $J$, integrality follows directly from
$\delta_1\equiv\delta_2$ modulo two. To prove it for
$H$, first put $K=(\delta_1+\delta_2)/2\in\mathbb Z[D]$.
The actual coefficient $G_1=3DK$ is zero modulo two,
so cancellation in $\mathbb F_2[D]$ gives $K\equiv0$;
therefore $H=K/2$ is integral. Effectivity follows
from their displayed positive combinations of the cells.
The $G_2$ parity gives $J\equiv D(D-1)$ modulo two.

Let $\sigma=q^2+q+q^{-1}+q^{-2}=D(D-3)$ denote
the absolute cell trace coefficient. The actual cell
bounds in (G028.2) are exactly
\[
 2H+\sigma\le J\le6H-\sigma                       \tag{G028.3}
\]
coefficientwise. In particular $H$ has coefficients
at least one at each of $q^{\pm1},q^{\pm2}$.
Write
\[
 m=H(4)\ge4,\qquad n=J(4)\ge2m+4.
\]
Here $m$ is an integer and $n$ is a positive even integer.
No degree or Laurent-support bound has been imposed.

\paragraph{A mass inequality removes every $m\ge6$.}
Let $v_i=V_i(4)$ and $a_i=A_i(4)$ in this paragraph;
these symbols are masses, not the ratio numerators of
NN-G026. The actual coefficients obey
\[
 a_1=96m-(n/3)v_1\ge8,\qquad
 a_2=16n-2mv_2\ge16.
\]
The exact total mass becomes
\[
 \begin{split}
 T={}&8mn+584m+80n+2720+
       \frac{5760m}{n}+\frac{160n}{m}\\
    &+\left(14-\frac{60}{n}\right)a_1
       +\left(14-\frac{10}{m}\right)a_2.
 \end{split}                                                        \tag{G028.4}
\]
Both coefficients on the last line are positive,
since $m\ge4$ and $n\ge12$. Therefore
\[
 T\ge L(m,n):=8mn+584m+80n+3056
       +\frac{5760m-480}{n}+\frac{160n-160}{m}.       \tag{G028.5}
\]
For $m\ge6$ and $n\ge2m+4$, this is minimized at
$(m,n)=(6,16)$. Here is a direct verification of that
global assertion, rather than a bounded numerical test.
The derivative with respect to $n$ is
\[
 8m+80+160/m-(5760m-480)/n^2.
\]
Its value at $n=2m+4$, multiplied by $m(m+2)^2$,
is
\[
 8\{m^2(m^2+14m-116)+135m+80\}>0
 \quad(m\ge6).
\]
It is increasing in $n$, so $L(m,n)\ge L(m,2m+4)$.
The latter expression is
\[
 16m^2+776m+6576-\frac{6000}{m+2}+\frac{480}{m},
\]
whose derivative
$32m+776+6000/(m+2)^2-480/m^2$ is positive for
$m\ge6$. Thus
\[
 T\ge L(6,16)=11138>11128.                          \tag{G028.6}
\]
The new attainer has total mass $11128$, so every
competing minimizer has $m\le5$.

If $m=4$, (G028.3) forces $H=\sigma$.
Then $[q^0]G_1=[q^0](6D\sigma)=12<16$, violating
its trace target. If $m=5$, the only possible addition
to those four mandatory coefficients is one central
unit, so
\[
 H=(1,1,1)=D^2-3D+1.                                \tag{G028.7}
\]
This proves a degree reduction from the score and
actual coefficient inequalities; it does not infer
a degree bound from mass alone.

\paragraph{Only four $J$ polynomials remain.}
By (G028.3) the support of $J$ is contained in that
of $H$, so write its symmetric coefficient list as
$(j_0,j_1,j_2)$. Its bounds and parity give
\[
 j_0\in\{2,4,6\},\qquad j_1,j_2\in\{3,5\}.
\]
Effectivity of $A_i=DG_i-(G_j/W)V_i$ forces
$\deg V_i\le2$: above that degree its highest term
would have negative sign, since $H,J,V_i$ have
positive leading coefficients. The coefficient of
$V_2$ at $q^2$ is even and at least two. The top
coefficient of $A_2=D^2J-2HV_2$ is therefore at most
$j_2-4$, forcing $j_2=5$; its $V_2$ coefficient is
then exactly two.

The central coefficient of $S=2DHJ$ equals
$8j_0+16j_1+60$ and must be at least $144$.
This leaves precisely
\[
 (j_0,j_1,j_2)=(2,5,5),(4,5,5),(6,3,5),(6,5,5).
 \tag{G028.8}
\]
Each $J$ in this list is primitive. Since
$A_1=6D^2H-(J/3)V_1$ is integral, Gauss's lemma
forces $V_1=3C$ with $C$ an integral effective
polynomial. The parity and target of $V_1$ give
$C_0\ge2$ even, $C_1\ge3$ odd, and $C_2$ even.
The leading inequality $6-5C_2\ge0$ forces
$C_2=0$. Thus $C=(x,y)$ with $x\ge2$ even and
$y\ge3$ odd.

\paragraph{Complete coefficient optimization of the two remaining factors.}
The symmetric list of $6D^2H$ is
$(96,90,66,30,6)$. Subtracting $JC$ and comparing
with $X_1=(2,1,1,1)$ gives exactly
\[
 \begin{gathered}
 y\in\{3,5\},\qquad
 5x+j_1y\le65,\\
 j_1x+(j_0+5)y\le89,\qquad
 j_0x+2j_1y\le94.                                  \end{gathered}\tag{G028.9}
\]
For fixed $y$, the largest even $x$ satisfying these
three bounds maximizes $C(4)=x+2y$.

Write $V_2=(z,t,2)$, where $z\ge6$ is even and
$t\ge3$ is odd. Expanding $D^2J-2HV_2$ and comparing
with $X_2=(2,3,2,1,1)$ gives exactly
\[
 \begin{gathered}
 t\le(j_1+15)/2,\qquad
 z+2t\le3j_0+4j_1,\\
 z+2t\le2j_0+(7j_1+13)/2,\qquad
 z+t\le j_0/2+2j_1+12.                             \end{gathered}\tag{G028.10}
\]
Here $t$ ranges only over $3,5,7,9$. For fixed $t$
take the largest allowed even $z$. These lists
include every coefficient condition: the remaining
top coefficients are $6$ and $1$, and parity is
already enforced. The total mass is strictly decreasing
in $v_1$ and $v_2$, with coefficients
$20-14n/3$ and $20-28m$, so maximizing these masses
independently is the exact minimizing rule.

The finite arithmetic is displayed in full by
\texttt{verify\_fixed\_bottom\_optimum.py}; its four
rows are
\[
\begin{array}{c|c|c|c|c|r}
 J&\max C(4)&\max V_2(4)&a_1&a_2&\text{score}\\\hline
 (2,5,5)&18&30&84&52&3720\\
 (4,5,5)&18&36&48&24&3452\\
 (6,3,5)&20&32&40&32&3352\\
 (6,5,5)&16&38&64&36&3728
\end{array}                                                         \tag{G028.11}
\]
There is only one minimizing $J$. In that row (G028.9)
forces the maximizing $C=(10,5)$, so $V_1=15D$.
Equation (G028.10) has exactly two maximizing choices,
\[
 V_2=(10,9,2)=2D^2+D-4,
 \quad\text{or}\quad
 V_2=(14,7,2)=2D^2-D+4.
\]
The first is the verified construction (G027.4).
For the second the actual $A_2$ list is
$(6,5,2,5,1)$; it has the same mass and satisfies
all required conditions. The other data are unchanged,
and the generic block construction again proves
associativity. This proves both attainment and the
complete two-parameter-set minimum in the class
(G028.1).\hfill$\square$

\paragraph{Geometric translation and the next exact gate.}
Fixing (G028.1) fixes three actual multiplicity objects
in the NN-G010 finite-kernel presentation. Among all
its effective graded sign-representation tables with
those fixed objects and the specified trace specialization,
the sum over multiplication positions and grading pieces
of half the dimension-minus-absolute-trace excess is
exactly $3352$. There are two minimizing character-valued
parameter sets in this presentation; no uniqueness of
categories or coherence data is asserted. This is an
optimization of chosen multiplication data, not an
intrinsic minimum for every categorification of the
abstract trace algebra.

The exact source dictionary remains SGA 3, I,
4.7.2--4.7.3 for effective grading characters and
archive cote 2, scan 17 for kernel composition.
The mass reduction respects actual coefficient
objects and never treats a rational auxiliary factor
as an effective object. A smaller closed-table score
must change at least one of $W,U_1,U_2$ in (G028.1);
other table shapes are a further separate possibility.
No canonicality, rank-four, PH0/PH1 or purity theorem
is inferred from this exact restricted optimum.
